\documentclass[11pt,letterpaper]{article}
\usepackage[margin=1in]{geometry}
\usepackage{amsmath,amssymb,booktabs,tabularx}
\setlength{\parindent}{0pt}
\setlength{\parskip}{5pt}
\setlength{\emergencystretch}{2em}
\raggedbottom
\renewcommand{\arraystretch}{1.15}
\title{Decision framework for a hybrid charging network}
\author{}
\date{}
\begin{document}
\maketitle


This framework describes an integrated operator managing Fixed Charging Stations (FCSs) and Mobile Charging Vehicles (MCVs) to maximize expected cumulative operational profit. MCVs can deliver charging energy to EV customers or export battery energy to the grid through an FCS connection.

The policy symbols below describe proposed learning components. The current simulator implements the system dynamics, observations, and rewards; it does not implement or train these policies.

\section*{1. Two-stage decision architecture}


Each physical period follows a pricing decision, customer choice, and a dispatch decision:

\[
s_t
\xrightarrow{\text{Stage 1: pricing}} p_t^M
\xrightarrow{\text{Customer choice}} D_t
\]


\[
\widetilde{s}_t
\xrightarrow{\text{Stage 2: dispatch}} a_t^D
\xrightarrow{\text{System transition}} (r_t,s_{t+1}).
\]


\subsection*{Stage 1: Pricing phase}


An abbreviated pricing state is

\[
s_t =
\left(
\{(l_{mt},e_{mt})\}_{m\in\mathcal{M}},
\lambda_t,p_t^F
\right).
\]


Here $l_{mt}$ and $e_{mt}$ are vehicle locations and stored energies, $\lambda_t$ denotes realized EV request counts by zone, and $p_t^F$ denotes FCS retail prices. The full implemented state also includes the period index and known grid buy and sell prices.

The operator quotes a price for each MCV using a proposed pricing policy:

\[
p_m\in[p_{\min},p_{\max}],
\qquad
p_t^M\sim\pi^P_{\theta_1}(\cdot\mid s_t).
\]


The bounds $p_{\min}$ and $p_{\max}$ correspond to \texttt{price\_min} and \texttt{price\_max}.

\subsection*{Customer choice}


Customers choose among available MCVs, FCSs, and an outside option after observing the quoted prices. Their random choices induce endogenous demand:

\[
D_t=D(p_t^M;\xi_t).
\]


Here $\xi_t$ represents customer-choice randomness.

\subsection*{Stage 2: Dispatch phase}


After customer choices, the post-choice state is

\[
\widetilde{s}_t=(s_t,p_t^M,D_t).
\]


An abbreviated local observation for MCV $m$ and its proposed dispatch policy are

\[
o_{mt}=(l_{mt},e_{mt},p_t^F,p_t^M,D_t),
\qquad
a^D_{mt}\sim\pi^D_{\theta_{2m}}(\cdot\mid o_{mt}).
\]


The full local observation also includes the period index and feasibility information. Dispatch advances time by one physical period and returns period profit $r_t$.

\section*{2. Exogenous request arrivals}


EV charging requests arrive independently across zones $z$ and periods $t$. Every customer requests the same energy quantity $q$, specified by \texttt{request\_kwh}.

The number of requests in zone $z$ at period $t$ follows

\[
N_{tz}\sim\mathrm{Binomial}
\left(
N_{\max},\frac{\mu_z h_t}{N_{\max}}
\right).
\]


\begin{itemize}
\setlength{\itemsep}{1pt}
\item $N_{\max}$: maximum requests per zone, \texttt{max\_requests\_per\_zone}.
\item $\mu_z$: baseline mean requests in zone $z$, \texttt{mean\_requests\_by\_zone[z]}.
\item $h_t$: arrival-profile multiplier, \texttt{arrival\_profile[t]}.
\end{itemize}


Validation requires $0\leq\mu_z h_t/N_{\max}\leq1$.

\section*{3. Endogenous customer choice model}


Customers evaluate available options using a multinomial logit model based on price and distance.

\subsection*{Utility equations}


For a customer in zone $z$, mobile option $m$ has utility

\[
u^M_{zm}=b_M-\alpha q p_m-\delta_M d(L_m,z).
\]


Fixed station $f$, located in zone $z_f$, has utility

\[
u^F_{zf}=b_F-\alpha q p^F_f-\delta_F d(z,z_f).
\]


The outside option has fixed utility $u_{\mathrm{out}}$.

\begin{itemize}
\setlength{\itemsep}{1pt}
\item $b_M,b_F$: base utility constants.
\item $\alpha$: sensitivity to the customer's total charging bill.
\item $p_m,p^F_f$: retail prices per delivered kWh.
\item $\delta_M,\delta_F$: distance-disutility coefficients.
\item $d(\cdot,\cdot)$: Euclidean distance in kilometers.
\item $L_m$: current zone of MCV $m$.
\end{itemize}


\subsection*{Choice probabilities}


The probability that a customer in zone $z$ chooses option $j$ is

\[
P_{zj}=
\frac{\exp(u_{zj})}
{\sum_{k\in\mathcal{A}_z}\exp(u_{zk})}.
\]


The set $\mathcal{A}_z$ contains the available options, including the outside option. An MCV is available only if its battery energy and remaining operating time allow at least one full request. An FCS is available if it lies within the access radius and its EV-service capacity permits at least one full request. Unavailable options have probability zero.

Availability does not guarantee that all customers who select an option will receive service. An FCS can attract more requests than it can fulfill. An MCV can attract requests in several zones but serve only one zone in the period. Unserved customers do not queue, retry another provider, or carry demand forward.

\section*{4. Dispatch actions and physical constraints}


Each MCV chooses one mode and one target zone per period.

\begin{center}
\begin{tabularx}{\textwidth}{@{}lX@{}}
\toprule
\textbf{Mode} & \textbf{Function} \\
\midrule
\texttt{WAIT} & Stay in the current zone without movement or energy exchange. \\
\texttt{SERVE} & Travel to a target zone and deliver energy to selected EV customers. \\
\texttt{REPOSITION} & Move to a zone without serving or exchanging energy. \\
\texttt{RECHARGE} & Travel to an FCS and purchase grid energy to charge the MCV battery. \\
\texttt{DISCHARGE} & Travel to an FCS and export battery energy to the grid. \\
\bottomrule
\end{tabularx}
\end{center}


\subsection*{Travel time and energy}


For distance $d$, the remaining operating time and travel energy consumption are

\[
h=\Delta-\frac{d}{v},
\qquad
e^{\mathrm{travel}}=\kappa d.
\]


Here $\Delta$ is period duration (one hour by default), $v$ is speed in km/h, and $\kappa$ is travel energy consumption in kWh/km. Travel must fit within the period and leave the battery reserve intact.

\subsection*{Serving capacity}


The maximum number of full requests MCV $m$ can serve in zone $z$ is

\[
k_{mz}=
\left\lfloor
\frac{
\min\left\{
\eta_M(E_m-e^{\mathrm{travel}}-E_{\min}),\,hP_M
\right\}_+
}{q}
\right\rfloor.
\]


Here $E_m$ is stored battery energy, $E_{\min}$ is the reserve, $\eta_M$ is service efficiency, and $P_M$ is delivered service power in kW. The notation $x_+=\max(x,0)$ clips negative capacity to zero. Actual service is limited by both this capacity and the customers who selected that vehicle in the target zone.

\section*{5. Station capacity and proportional scaling}


Each FCS $f$ has a per-period fleet exchange budget $C_f$, specified by \texttt{fcs\_fleet\_exchange\_capacity\_kwh}. This budget is separate from the FCS's EV-service capacity.

Let $x_m$ be an MCV's requested grid-side charging or export energy at station $f$. For positive total requests, the allocated amount is

\[
\widehat{x}_m=
x_m\min\left\{
1,\frac{C_f}{\sum_{j\text{ at }f}x_j}
\right\}.
\]


If total requests are zero, the scaling factor is defined as one and all allocations are zero. Charging and discharging share a gross energy budget without netting.

\section*{6. Energy conservation and state update}


Stored battery energy at the end of the period obeys

\[
E'_m=
E_m-e_m^{\mathrm{travel}}
+\eta_C C_m
-\frac{Q_m}{\eta_M}
-\frac{X_m}{\eta_D}.
\]


\begin{itemize}
\setlength{\itemsep}{1pt}
\item $C_m$: allocated grid-input energy used to charge the MCV; $\eta_C$ is charging efficiency.
\item $Q_m$: energy delivered to EV customers; $\eta_M$ is service efficiency.
\item $X_m$: energy exported to the grid; $\eta_D$ is discharging efficiency.
\end{itemize}


The feasible final inventory satisfies

\[
E_{\min}\leq E'_m\leq E_{\max},
\]


where $E_{\max}$ is the configured battery capacity. Energy quantities are measured in kWh; power quantities are measured in kW.

\section*{7. Financial objective and reward dynamics}


Period profit includes retail revenue, grid transactions, operating costs, unmet-demand penalties, and final-period salvage value:

\[
\begin{aligned}
r_t={}&
\sum_m p_m Q_m+\sum_f p^F_f Q^F_f
+g_t^{\mathrm{sell}}\sum_m X_m\\
&-g_t^{\mathrm{buy}}
\left(\sum_m C_m+\sum_f\frac{Q^F_f}{\eta_F}\right)
-c_{\mathrm{travel}}\sum_m d_m\\
&-c_{\mathrm{op}}|\mathcal{M}|
-c_{\mathrm{unmet}}qU_t
+\mathbf{1}_{\{t=T-1\}}s\sum_m E'_m.
\end{aligned}
\]


\subsection*{Revenue and cost terms}


\begin{itemize}
\setlength{\itemsep}{1pt}
\item $\sum_m p_mQ_m+\sum_f p^F_fQ^F_f$: retail revenue from energy delivered to EVs.
\item $g_t^{\mathrm{sell}}\sum_mX_m$: grid-export revenue.
\item $g_t^{\mathrm{buy}}(\sum_mC_m+\sum_fQ^F_f/\eta_F)$: electricity purchase cost, accounting for FCS service efficiency $\eta_F$.
\item $c_{\mathrm{travel}}\sum_md_m$: travel cost.
\item $c_{\mathrm{op}}|\mathcal{M}|$: operating cost for every MCV in every period, including vehicles that wait.
\item $c_{\mathrm{unmet}}qU_t$: penalty for customers who selected an operator service but were not served; outside-option customers are excluded.
\item $\mathbf{1}_{\{t=T-1\}}s\sum_mE'_m$: terminal inventory value at $s$ dollars per stored kWh.
\end{itemize}


Here $Q^F_f$ is energy delivered to EVs by FCS $f$, $d_m$ is MCV travel distance, and $g_t^{\mathrm{buy}}$ and $g_t^{\mathrm{sell}}$ are the known grid buy and sell prices.

\subsection*{Long-term optimization objective}


The intended learning objective is expected cumulative discounted profit:

\[
\max_{\pi^P,\pi^D}
\mathbb{E}\left[
\sum_{t=0}^{T-1}\gamma^t r_t
\right],
\qquad 0\leq\gamma\leq1.
\]


The parameter $\gamma$ corresponds to \texttt{discount\_per\_period}. Discounting applies once per physical period. The pricing transition has continuation discount one; a nonterminal dispatch transition has continuation discount $\gamma$, and the terminal dispatch transition has continuation discount zero.

\end{document}
